\section{Conclusion}
\label{sec:conclusion}
Energy-matching samplers with self-normalized targets, such as iDEM, regress onto Monte Carlo targets whose population minimizer ignores the buffer, so the target distribution is not a fixed point of their training. We showed how to measure this inconsistency from reference samples alone, found it to be large on the standard molecular benchmarks, and removed it by putting back the one exact posterior sample every training pair already contains. The fix costs one line of code and one energy evaluation, and it comes with unbiasedness, variance and contraction guarantees.
Our experiments also separate two effects that are usually measured together. A consistent target fixes the within-mode distribution and the learned score; it does not discover modes, and on our molecular runs it did not by itself give better samplers. Successors with bootstrapped targets, and the original implementations of iEFM and CVSI, can be checked the same way; we tested only our implementations of their target forms. We suggest that new energy-matching objectives report the consistency test of \cref{prop:diag}, which is cheap, and that consistent targets be paired with exploration that does not rely on biased estimators (tempering, annealed importance sampling, parallel chains across noise levels).
\paragraph{Limitations.}
Our targets have at most 39 dimensions; training budgets are short ($3000$--$8000$ steps), with 3 seeds and no per-method tuning, and several comparisons are within one standard deviation. LJ-13 uses a smaller network and sampler grid. The molecular test and oracle runs are relative to MALA references, which may miss isolated basins. Larger budgets may change the ranking on mode-level and molecular metrics. We compare only with iDEM in training.
